% !TeX root = main.tex
% LTeX: language=en-US

\begin{frame}
    \titlepage{}
\end{frame}

\begin{frame}{Table of Contents}
    \tableofcontents
\end{frame}

%==========================================================
\section{Introduction}

\begin{frame}{Context}
    \begin{BlockBlue}{\emph{maybe} Kent Beck citation}
        \begin{enumerate}
            \item \makebox[7em][l]{\emph{Make it work}}\only<2->{\( \leftarrow \) done!}
            \item \makebox[7em][l]{\emph{Make it right}}\only<2->{\( \leftarrow \) done!}
            \item \only<1>{\emph{Make it fast}}\only<2->{\makebox[7em][l]{\textbf{\emph{Make it fast}}}\textbf{\( \leftarrow \) to do!}}
        \end{enumerate}
    \end{BlockBlue}

    \bigskip

    \visible<3->{%
        \begin{BlockBlue}{Donald Knuth}
            \itshape
            \visible<4>{We should forget about small efficiencies, say about 97\% of the time:}\par
            Premature optimization is the root of all evil.
            \visible<4>{Yet we should not pass up our opportunities in that critical 3\%.}
        \end{BlockBlue}
    }
\end{frame}


\begin{frame}{Slow code, fast Julia}
    \begin{BlockBlue}{Why is code slow? By potential gain}
        \begin{itemize}%\setlength{\itemsep}{0pt}\setlength{\topsep}{0pt}
            \item \textbf{Language / compiler}: in Julia, type stability and allocations
            \item \textbf{Memory / architecture}: caches, access order, SIMD (4--8 operations at once)
            \item \textbf{Algorithm}: complexity, redundant computations, existing library
            \item \textbf{Parallelism}: threads, MPI, GPU (other sessions)
        \end{itemize}
    \end{BlockBlue}

    \pause\smallskip

    \begin{BlockGreen}{Why can Julia be fast? Compilation for concrete types}
        \begin{itemize}%\setlength{\itemsep}{0pt}\setlength{\topsep}{0pt}
            \item Each function is compiled (LLVM) \textbf{for the concrete types of its arguments}
            \item If the compiler knows all the types \( \to \) machine code comparable to C/Fortran
            \item Otherwise \( \to \) dynamic dispatch, ``boxed'' values, allocations: slower
        \end{itemize}
    \end{BlockGreen}
\end{frame}


\begin{frame}{The method}
    \begin{enumerate}
        \item \textbf{Baseline}: measure + correctness test
        \item \textbf{Profile}: where is the time spent?
        \item \textbf{Hypothesis}: why here?
        \item \textbf{One} change
        \item \textbf{Re-measure} + check the result
        \item \textbf{keep or revert}: start again at 2.
    \end{enumerate}

    \bigskip

    \alert{\textbf{Never optimize before measuring.}}
\end{frame}

%==========================================================
\section{Benchmark}

\subsection{Exercise 1: Measuring correctly}

\begin{frame}{Exercise 1: Measuring correctly}
    \texttt{exercises/01\_measure.ipynb}

    \pause\bigskip

    \begin{BlockGreen}{What we learned}
        \begin{itemize}
            \item 1st call = JIT compilation: ignore it; a single run is noisy
            \item \textbf{Performance-critical code goes inside functions}
            \item \texttt{@btime} / \texttt{@benchmark}, and \texttt{\$} to interpolate globals
            \item Mutating function: \texttt{setup=(\ldots) evals=1}
            \item Floats: compare with \texttt{\(\approx\)}, never \texttt{==}
        \end{itemize}
    \end{BlockGreen}
\end{frame}

%==========================================================
\subsection{Exercise 2: Type stability}

\begin{frame}{Exercise 2: Type stability}
    \texttt{exercises/02\_types.ipynb}

    \pause\bigskip

    \begin{BlockGreen}{What we learned}
        \begin{itemize}
            \item Diagnose with \texttt{@code\_warntype} (red \texttt{Any}) and \texttt{@report\_opt} (JET)
            \item Non-\texttt{const} global \( \to \) \texttt{const} or argument (\( \times 40 \))
            \item Field of abstract type \( \to \) parametric \texttt{struct} (\( \times 70 \))
            \item \texttt{Dict\{String,Any\}} \( \to \) function barrier (\( \times 150 \))
            \item Type instability often shows up as unexpected allocations
        \end{itemize}
    \end{BlockGreen}
\end{frame}

%==========================================================
\subsection{Exercise 3: Allocations and GC}

\begin{frame}{Exercise 3: Allocations and GC}
    \texttt{exercises/03\_allocations.ipynb}

    \pause\bigskip

    \begin{BlockGreen}{What we learned}
        \begin{itemize}
            \item \texttt{@benchmark}: the GC shows up in the mean time, not in the min
            \item Temporaries: fusion \texttt{.}, \texttt{sum(f, x)}, buffer \texttt{f!}
            \item Slices copy: \texttt{@views}; growing arrays: preallocate, \texttt{sizehint!}
            \item Small fixed-size vectors: \texttt{SVector}, 0 allocations
            \item \alert{Goal: no allocations in critical loops}
        \end{itemize}
    \end{BlockGreen}
\end{frame}


%==========================================================
\section{Profiling}

\subsection{Exercise 4: Profiling}

\begin{frame}{Exercise 4: Profiling}
    \texttt{exercises/04\_profiling.ipynb}

    \pause\bigskip

    \begin{BlockGreen}{What we learned}
        \begin{itemize}
            \item Run once before profiling; profile a long enough run
            \item Flame graph: width = time; \Rd{\textbf{red}} = dispatch; \Or{\textbf{yellow}} = GC
            \item \texttt{@profview\_allocs}: where the allocations come from
            \item \textbf{After each fix, profile again}: the hot spot moves
            \item 1.27 s \( \to \) 10.5 ms (\( \times 120 \))
        \end{itemize}
    \end{BlockGreen}
\end{frame}

%==========================================================
\subsection{Exercise 5: Memory and cache}

\begin{frame}{Exercise 5: Memory and cache}
    \texttt{exercises/05\_memory.ipynb}

    \pause\bigskip

    \begin{BlockGreen}{What we learned}
        \begin{itemize}
            \item Column-major: inner loop over the \textbf{first} index
            \item The throughput drops in steps when the data leave L2, then L3
            \item Matvec (0.25 flop/byte): memory-bound, \texttt{mul!} \( \approx \) a good loop
            \item Matmul (\( n/12 \) flop/byte): compute-bound, BLAS \( \times 5 \) (\( \times 18 \) with 4 threads)
            \item \alert{Before writing a kernel, check the libraries}
        \end{itemize}
    \end{BlockGreen}
\end{frame}

%==========================================================
\section{Case study}

\subsection{Exercise 6: Case study}

\begin{frame}{Exercise 6: 2D heat equation}
    \texttt{exercises/06\_case\_study.ipynb}

    \pause\bigskip

    \begin{BlockGreen}{What we learned}
        \begin{itemize}
            \item Type stability first: most of the gain (\( \times 48 \))
            \item Then allocations (\( \times 56 \)), then loop order (\( \times 108 \))
            \item \texttt{@inbounds @simd}: \( \times 260 \), once the code is clean
            \item Threads bring little on a small grid
            \item \textbf{One change at a time}, check after each one
        \end{itemize}
    \end{BlockGreen}
\end{frame}

%==========================================================
\section{Wrap-up}

\begin{frame}{Checklist}
    \begin{enumerate}
        \item \textbf{Algorithm}: complexity? redundancy? library?
        \item \textbf{Measure}: \texttt{@btime f(\$x)} + correctness test
        \item \textbf{Types}: \texttt{@code\_warntype}, \texttt{@report\_opt}, \texttt{const}, parametric \texttt{struct}s, barriers
        \item \textbf{Allocations}: \% of GC, \texttt{@views}, fusion, buffers, \texttt{StaticArrays}
        \item \textbf{Profile}: \texttt{@profview}, \texttt{@profview\_allocs}, profile again
        \item \textbf{Memory}: loop order, contiguity; memory- or compute-bound?
        \item \textbf{Only then}: SIMD, threads, MPI, GPU
    \end{enumerate}
\end{frame}

\begin{frame}{Later this week}
    \begin{BlockBlue}{Related sessions}
        \begin{itemize}
            \item \textbf{Shared-memory parallelism}: Threads.jl, DataFlowTasks.jl
            \item \textbf{GPU}: a different execution model (transfers, massive parallelism)
            \item \textbf{Multiple dispatch and generic typing}: parametric types, generic code
            \item \textbf{MPI}: distributed-memory parallelism
        \end{itemize}
    \end{BlockBlue}

    \bigskip

    \alert{\textbf{Parallelize code that is already fast on one core.}}
\end{frame}

\begin{frame}{Resources}
    \begin{itemize}
        \item Links:
              \begin{itemize}
                  \item \href{https://docs.julialang.org/en/v1/manual/profile/}{\Pu{Profiling}}, Julia manual
                  \item \href{https://docs.julialang.org/en/v1/manual/performance-tips/}{\Pu{Performance Tips}}, Julia manual
                  \item \href{https://book.sciml.ai/notes/18-Code_Profiling_and_Optimization/}{\Pu{Code Profiling and Optimization}}, SciML book
              \end{itemize}

        \item Standard libraries:
              \begin{itemize}
                  \item Profile, LinearAlgebra, Test
              \end{itemize}

        \item Third-party packages:
              \begin{itemize}
                  \item BenchmarkTools.jl, JET.jl, StaticArrays.jl, ProfileCanvas.jl
              \end{itemize}
    \end{itemize}

    \bigskip

    \centering
    \alert{\textbf{Measure \( \to \) Understand \( \to \) Optimize \( \to \) Verify}}
\end{frame}
